\documentclass[12pt]{article}
\usepackage[utf8]{inputenc}
\usepackage[T1]{fontenc}
\usepackage{hyperref}
\title{Sonaga Language Resources for K-12 Curriculum Development}
\author{Compiled by Claude Sonnet 5.5 (claude-sonnet-5-5) for Andr\'as Kornai}
\date{}
\begin{document}
\maketitle

\section{Introduction}
Sonaga (ISO 639-3: ysg) is a Loloish language of the Sino-Tibetan family, spoken in Heqing County, Yunnan, China. Wikipedia places its speakers mainly in Dongdeng of Xinfeng Village, Caohai Township, with some in Anle Village. It reports about 2,000 native speakers (2009); Joshua Project gives 2,200. Chinese sources call it 锁内嘎话, and it is also known as Xinfeng Yi. The searches found no educational materials for Sonaga. The sources below are general reference pages only.

\section{Primary Educational Resources}
No primers, textbooks or school programmes for Sonaga were found in the searches. The language is not known to be taught in schools.

\section{Linguistic and Reference Materials}
\begin{itemize}
\item Wikipedia, ``Sonaga language'': overview of classification, location and speaker numbers. \url{https://en.wikipedia.org/wiki/Sonaga_language}
\item Numeral list for Sonaga, Yunnan, collected by Andy Castro in 2009. \url{https://lingweb.eva.mpg.de/channumerals/Sonaga.htm}
\item The Wikipedia summary cites a 2010 SIL survey of Kua-nsi and related Yi varieties by Castro, Crook and Flaming. I did not find the survey itself in the search results.
\end{itemize}

\section{Religious and Community Resources}
\begin{itemize}
\item Joshua Project profile, ``Sonaga in China''. This is a missions demographic profile, not a community-produced resource. \url{https://joshuaproject.net/people_groups/20590}
\item Joshua Project language page for ysg. \url{https://joshuaproject.net/languages/ysg}
\end{itemize}

\section{Digital and Audio Resources}
No audio, video, corpus or software resources were found.

\section{Curriculum Development Context}
Sonaga has very few speakers and almost no documentation beyond a survey and a numeral list. Any curriculum would first need community-led documentation: a word list, an orthography decision, and recorded speech. The searches show no existing teaching materials to build on.
\end{document}